\documentclass[10pt,a4paper]{amsart}
\usepackage{amsmath,amssymb,amsthm,graphicx}
\usepackage{verbatim}
\newenvironment{reference}{\comment}{\endcomment}
%\newenvironment{reference}{}{}
\newenvironment{slogan}{\comment}{\endcomment}
\newenvironment{history}{\comment}{\endcomment}
% For commutative diagrams we use Xy-pic
\usepackage[all]{xy}

% We use 2cell for 2-commutative diagrams.
\xyoption{2cell}
% Theorem environments.
%
\theoremstyle{plain}
\newtheorem{theorem}[subsection]{Theorem}
\newtheorem{proposition}[subsection]{Proposition}
\newtheorem{lemma}[subsection]{Lemma}

\theoremstyle{definition}
\newtheorem{definition}[subsection]{Definition}
\newtheorem{example}[subsection]{Example}
\newtheorem{exercise}[subsection]{Exercise}
\newtheorem{situation}[subsection]{Situation}

\theoremstyle{remark}
\newtheorem{remark}[subsection]{Remark}
\newtheorem{remarks}[subsection]{Remarks}

\numberwithin{equation}{subsection}

% Macros
%
\def\lim{\mathop{\mathrm{lim}}\nolimits}
\def\colim{\mathop{\mathrm{colim}}\nolimits}
\def\Spec{\mathop{\mathrm{Spec}}}
\def\Hom{\mathop{\mathrm{Hom}}\nolimits}
\def\Ext{\mathop{\mathrm{Ext}}\nolimits}
\def\SheafHom{\mathop{\mathcal{H}\!\mathit{om}}\nolimits}
\def\SheafExt{\mathop{\mathcal{E}\!\mathit{xt}}\nolimits}
\def\Sch{\mathit{Sch}}
\def\Mor{\mathop{\mathrm{Mor}}\nolimits}
\def\Ob{\mathop{\mathrm{Ob}}\nolimits}
\def\Sh{\mathop{\mathit{Sh}}\nolimits}
\def\NL{\mathop{N\!L}\nolimits}
\def\CH{\mathop{\mathrm{CH}}\nolimits}
\def\proetale{{pro\text{-}\acute{e}tale}}
\def\etale{{\acute{e}tale}}
\def\QCoh{\mathit{QCoh}}
\def\Ker{\mathop{\mathrm{Ker}}}
\def\Im{\mathop{\mathrm{Im}}}
\def\Coker{\mathop{\mathrm{Coker}}}
\def\Coim{\mathop{\mathrm{Coim}}}

% Boxtimes
%
\DeclareMathSymbol{\boxtimes}{\mathbin}{AMSa}{"02}

%
% Macros for moduli stacks/spaces
%
\def\QCohstack{\mathcal{QC}\!\mathit{oh}}
\def\Cohstack{\mathcal{C}\!\mathit{oh}}
\def\Spacesstack{\mathcal{S}\!\mathit{paces}}
\def\Quotfunctor{\mathrm{Quot}}
\def\Hilbfunctor{\mathrm{Hilb}}
\def\Curvesstack{\mathcal{C}\!\mathit{urves}}
\def\Polarizedstack{\mathcal{P}\!\mathit{olarized}}
\def\Complexesstack{\mathcal{C}\!\mathit{omplexes}}
% \Pic is the operator that assigns to X its picard group, usage \Pic(X)
% \Picardstack_{X/B} denotes the Picard stack of X over B
% \Picardfunctor_{X/B} denotes the Picard functor of X over B
\def\Pic{\mathop{\mathrm{Pic}}\nolimits}
\def\Picardstack{\mathcal{P}\!\mathit{ic}}
\def\Picardfunctor{\mathrm{Pic}}
\def\Deformationcategory{\mathcal{D}\!\mathit{ef}}
\usepackage{cleveref}
\crefname{originallemma}{original lemma}{original lemmas}
\Crefname{originallemma}{Original lemma}{Original lemmas}
\makeatletter
\newlabel{lemma-annihilated}{{lemma-annihilated}{source}{}{sourcepointer-1}{}}
\newlabel{lemma-annihilated@cref}{{[originallemma][0][]lemma-annihilated}{[1][0][]source}{}{}{}}
\makeatother
\crefname{originallemma}{original lemma}{original lemmas}
\Crefname{originallemma}{Original lemma}{Original lemmas}
\makeatletter
\newlabel{lemma-bound-q-and-d}{{lemma-bound-q-and-d}{source}{}{sourcepointer-2}{}}
\newlabel{lemma-bound-q-and-d@cref}{{[originallemma][0][]lemma-bound-q-and-d}{[1][0][]source}{}{}{}}
\makeatother
\crefname{originallemma}{original lemma}{original lemmas}
\Crefname{originallemma}{Original lemma}{Original lemmas}
\makeatletter
\newlabel{lemma-vanishing-coh-almost-projective}{{lemma-vanishing-coh-almost-projective}{source}{}{sourcepointer-3}{}}
\newlabel{lemma-vanishing-coh-almost-projective@cref}{{[originallemma][0][]lemma-vanishing-coh-almost-projective}{[1][0][]source}{}{}{}}
\makeatother
\crefname{originalreference}{original reference}{original references}
\Crefname{originalreference}{Original reference}{Original references}
\makeatletter
\newlabel{more-algebra-lemma-dual-near-projective}{{more-algebra-lemma-dual-near-projective}{source}{}{sourcepointer-4}{}}
\newlabel{more-algebra-lemma-dual-near-projective@cref}{{[originalreference][0][]more-algebra-lemma-dual-near-projective}{[1][0][]source}{}{}{}}
\makeatother
\crefname{originalremark}{original remark}{original remarks}
\Crefname{originalremark}{Original remark}{Original remarks}
\makeatletter
\newlabel{remark-duals}{{remark-duals}{source}{}{sourcepointer-5}{}}
\newlabel{remark-duals@cref}{{[originalremark][0][]remark-duals}{[1][0][]source}{}{}{}}
\makeatother
\title{Factoring a map from a uniform module through a sufficiently negative shift of the ideal powers}
\author{The Stacks Project Authors}
\date{stacks-project:local-cohomology.tex:local-snapshot-20261012; GNU FDL 1.2 or later}
\begin{document}
\maketitle
\noindent\textit{Editorial note.} All delivered mathematical text is the author's own, copied from the pinned Stacks Project chapter source sources/stacks-local-cohomology/pinned-original-source.tex; no statement or proof step was rewritten, reordered or removed.\par\medskip
\noindent\textit{Editorial note.} Preamble provenance: the Stacks Project keeps one single shared preamble for every chapter; this card takes it directly from the staged sources/stacks-local-cohomology/preamble.tex as three verbatim spans (lines 8-12, 14-18 and 154-217), so no derived file was written and nothing was added or redefined. The wrapper class and the base packages are supplied by the assembler.\par\medskip
\noindent\textit{Editorial note.} Bibliography provenance: neither the counted statement nor its proof contains a citation (this lane rescans the whole span before writing the card), so the card carries no bibliography and no entry was reconstructed.\par\medskip
\noindent\textit{Editorial note.} Selection provenance: this card was opened under the file-based control contract (CONTROL.md), with the mechanical pre-card gate run on its statement and proof interval before the card was written.\par\medskip
\noindent\textit{Allowed standard prerequisites (editorial).}\begin{itemize}
\item The notion of an (A, n, c)-module and the compatible system of canonical maps from the pullback of the ideal powers to the structure sheaves of the closed subschemes. Set up in the author's section on a bit of uniformity, where the counted statement occurs, and used at the start of the counted proof. Fixes the hypothesis on the module and the maps along which the given homomorphism is pulled back.
\item Duality for the module and the description of maps out of a near-projective module. Both are cross-chapter pointers invoked inside the counted proof. Convert the pullback of the given map into a map of dual modules, which is where the shift in the ideal power comes from.
\end{itemize}\medskip
\ifdefined\hypertarget\hypertarget{sourcepointer-1}{}\fi
\noindent\textit{Source reference.} Original lemma lemma-annihilated: sources/stacks-local-cohomology/pinned-original-source.tex, source lines 5241--5241. Navigation pointer only: the counted text refers to the Stacks item lemma-annihilated, which belongs to another chapter of the project and is not part of this card. The pointer carries the target's own label as its number because this chapter's staged snapshot does not contain the other chapter's tag table; its content is not reproduced here.\par\medskip
\ifdefined\hypertarget\hypertarget{sourcepointer-2}{}\fi
\noindent\textit{Source reference.} Original lemma lemma-bound-q-and-d: sources/stacks-local-cohomology/pinned-original-source.tex, source lines 5200--5200. Navigation pointer only: the counted text refers to the Stacks item lemma-bound-q-and-d, which belongs to another chapter of the project and is not part of this card. The pointer carries the target's own label as its number because this chapter's staged snapshot does not contain the other chapter's tag table; its content is not reproduced here.\par\medskip
\ifdefined\hypertarget\hypertarget{sourcepointer-3}{}\fi
\noindent\textit{Source reference.} Original lemma lemma-vanishing-coh-almost-projective: sources/stacks-local-cohomology/pinned-original-source.tex, source lines 5373--5373. Navigation pointer only: the counted text refers to the Stacks item lemma-vanishing-coh-almost-projective, which belongs to another chapter of the project and is not part of this card. The pointer carries the target's own label as its number because this chapter's staged snapshot does not contain the other chapter's tag table; its content is not reproduced here.\par\medskip
\ifdefined\hypertarget\hypertarget{sourcepointer-4}{}\fi
\noindent\textit{Source reference.} Original reference more-algebra-lemma-dual-near-projective: sources/stacks-local-cohomology/pinned-original-source.tex, source lines 5375--5375. Navigation pointer only: the counted text refers to the Stacks item more-algebra-lemma-dual-near-projective, which belongs to another chapter of the project and is not part of this card. The pointer carries the target's own label as its number because this chapter's staged snapshot does not contain the other chapter's tag table; its content is not reproduced here.\par\medskip
\ifdefined\hypertarget\hypertarget{sourcepointer-5}{}\fi
\noindent\textit{Source reference.} Original remark remark-duals: sources/stacks-local-cohomology/pinned-original-source.tex, source lines 5361--5361. Navigation pointer only: the counted text refers to the Stacks item remark-duals, which belongs to another chapter of the project and is not part of this card. The pointer carries the target's own label as its number because this chapter's staged snapshot does not contain the other chapter's tag table; its content is not reproduced here.\par\medskip
\begin{lemma}
\label{lemma-factor-hom}
With $q_0 = q(S)$ and $d = d(S)$ as above, let $M$ be an $(A, n, c)$-module
and let $\varphi : M \to I^n/I^{2n}$ be an $A$-linear map. Assume
$n \geq \max(q_0 + (1 + d)c, (2 + d)c)$ and if $d = 0$ assume
$n \geq q_0 + 2c$. Then the composition
$$
M \xrightarrow{\varphi} I^n/I^{2n} \to
I^{n - (1 + d)c}/I^{2n - (1 + d)c}
$$
is of the form $\sum a_i \psi_i$ with $a_i \in I^c$ and
$\psi_i : M \to I^{n - (2 + d)c}/I^{2n - (2 + d)c}$.
\end{lemma}

\begin{proof}
The case $d > 1$. Since we have a compatible system of maps
$p^*(I^q) \to \mathcal{O}_X(q)$ for $q \geq 0$ there are canonical maps
$p^*(I^q/I^{q + \nu}) \to \mathcal{O}_{Y_\nu}(q)$ for $\nu \geq 0$.
Using this and pulling back $\varphi$ we obtain a map
$$
\chi : p^*M \longrightarrow \mathcal{O}_{Y_n}(n)
$$
such that the composition
$M \to H^0(X, p^*M) \to H^0(X, \mathcal{O}_{Y_n}(n))$
is the given homomorphism $\varphi$ combined with the
map $I^n/I^{2n} \to H^0(X, \mathcal{O}_{Y_n}(n))$.
Since $\mathcal{O}_{Y_n}(n)$ is invertible on $Y_n$ the
linear map $\chi$ determines a section
$$
\sigma \in \Gamma(X, (p^*M)^\vee(n))
$$
with notation as in Remark \ref{remark-duals}.
The discussion in Remark \ref{remark-duals} shows
the cokernel and kernel of $can : p^*(M^\vee) \to (p^*M)^\vee$
are scheme theoretically supported on $Y_c$. By
Lemma \ref{lemma-annihilated} the map
$(p^*M)^\vee(n) \to (p^*M)^\vee(n - 2c)$ factors
through $p^*(M^\vee)(n - 2c)$; small detail omitted.
Hence the image of $\sigma$ in $\Gamma(X, (p^*M)^\vee(n - 2c))$
comes from an element
$$
\sigma' \in \Gamma(X, p^*(M^\vee)(n - 2c))
$$
By Lemma \ref{lemma-vanishing-coh-almost-projective} part (5),
the fact that $M^\vee$ is an $(A, n, c)$-module by
More on Algebra, Lemma \ref{more-algebra-lemma-dual-near-projective},
and the fact that $n \geq q_0 + (1 + d)c$ so $n - 2c \geq q_0 + (d - 1)c$
we see that the image of $\sigma'$ in $H^0(X, p^*M^\vee(n - (1 + d)c))$
is the image of an element $\tau$ in $I^{n - (1 + d)c}M^\vee$.
Write $\tau = \sum a_i \tau_i$ with $\tau_i \in I^{n - (2 + d)c}M^\vee$;
this makes sense as $n - (2 + d)c \geq 0$.
Then $\tau_i$ determines a homomorphism of modules
$\psi_i : M \to I^{n - (2 + d)c}/I^{2n - (2 + d)c}$
using the evaluation map $M \otimes M^\vee \to A/I^n$.

\medskip\noindent
Let us prove that this works\footnote{We hope some reader will suggest
a less dirty proof of this fact.}. Pick $z \in M$ and let us show that
$\varphi(z)$ and $\sum a_i \psi_i(z)$ have the same image in
$I^{n - (1 + d)c}/I^{2n - (1 + d)c}$.
First, the element $z$ determines a map
$p^*z : \mathcal{O}_{Y_n} \to p^*M$ whose composition with
$\chi$ is equal to the map $\mathcal{O}_{Y_n} \to \mathcal{O}_{Y_n}(n)$
corresponding to $\varphi(z)$ via the map
$I^n/I^{2n} \to \Gamma(\mathcal{O}_{Y_n}(n))$.
Next $z$ and $p^*z$ determine evaluation maps
$e_z : M^\vee \to A/I^n$ and $e_{p^*z} : (p^*M)^\vee \to \mathcal{O}_{Y_n}$.
Since $\chi(p^*z)$ is the section corresponding to $\varphi(z)$
we see that $e_{p^*z}(\sigma)$ is the section corresponding to $\varphi(z)$.
Here and below we abuse notation: for a map
$a : \mathcal{F} \to \mathcal{G}$ of modules on $X$
we also denote $a : \mathcal{F}(t) \to \mathcal{F}(t)$ the corresponding
map of twisted modules. The diagram
$$
\xymatrix{
p^*(M^\vee) \ar[d]_{can} \ar[r]_{p^*e_z} & \mathcal{O}_{Y_n} \ar@{=}[d] \\
(p^*M)^\vee \ar[r]^{e_{p^*z}} & \mathcal{O}_{Y_n}
}
$$
commutes by functoriality of the construction $can$. Hence
$(p^*e_z)(\sigma')$ in $\Gamma(Y_n, \mathcal{O}_{Y_n}(n - 2c))$
is the section corresponding to the image of $\varphi(z)$
in $I^{n - 2c}/I^{2n - 2c}$.
The next step is that $\sigma'$ maps to the image
of $\sum a_i \tau_i$ in $H^0(X, p^*M^\vee(n - (1 + d)c))$.
This implies that $(p^*e_z)(\sum a_i \tau_i) = \sum a_i p^*e_z(\tau_i)$
in $\Gamma(Y_n, \mathcal{O}_{Y_n}(n - (1 + d)c))$ is the section corresponding
to the image of $\varphi(z)$ in $I^{n - (1 + d)c}/I^{2n - (1 + d)c}$.
Recall that $\psi_i$ is defined from $\tau_i$ using an evaluation
map. Hence if we denote
$$
\chi_i : p^*M \longrightarrow \mathcal{O}_{Y_n}(n - (2 + d)c)
$$
the map we get from $\psi_i$, then we see by the same reasoning
as above that the section corresponding to $\psi_i(z)$ is
$\chi_i(p^*z) = e_{p^*z}(\chi_i) = p^*e_z(\tau_i)$. Hence we conclude that
the image of $\varphi(z)$ in $\Gamma(Y_n, \mathcal{O}_{Y_n}(n - (1 + d)c))$
is equal to the image of $\sum a_i\psi_i(z)$.
Since $n - (1 + d)c \geq q_0$ we have
$\Gamma(Y_n, \mathcal{O}_{Y_n}(n - (1 + d)c)) =
I^{n - (1 + d)c}/I^{2n - (1 + d)c}$ by Lemma \ref{lemma-bound-q-and-d} and
we conclude the desired compatibility is true.

\medskip\noindent
The case $d = 1$. Here we argue as above that we get
$$
\chi : p^*M \longrightarrow \mathcal{O}_{Y_n}(n),\quad
\sigma \in \Gamma(X, (p^*M)^\vee(n)),\quad
\sigma' \in \Gamma(X, p^*(M^\vee)(n - 2c)),
$$
and then we use
Lemma \ref{lemma-vanishing-coh-almost-projective} part (4)
to see that $\sigma'$ is the image of some element
$\tau \in I^{n - 2c}M^\vee$. The rest of the argument is the same.

\medskip\noindent
The case $d = 0$. Argument is exactly the same as in the
case $d = 1$.
\end{proof}

\end{document}
